\chapter{Strategi, Model Operasi, dan Visi Arsitektur}

\section*{Tujuan Pembelajaran}
\begin{itemize}
  \item Menurunkan pemicu, tujuan, dan \textit{outcome} organisasi menjadi kebutuhan arsitektur yang dapat ditelusuri.
  \item Menentukan model operasi organisasi berdasarkan tingkat integrasi dan standardisasi proses bisnisnya.
  \item Merumuskan visi arsitektur dan prinsip arsitektur lengkap dengan alasan dan implikasinya.
\end{itemize}

\section{Pendahuluan}

% TODO
% Cakupan: strategi direksi yang tidak pernah sampai ke keputusan sistem; proyek TI berjalan dan
%   selesai, tetapi tujuan strategis tidak tercapai.

\section{Dari Strategi ke Kebutuhan Arsitektur}

% TODO
% Cakupan: elemen motivasi (driver, assessment, goal, outcome, principle, requirement); rantai
%   keterlacakan dari strategi ke kebutuhan; outcome yang terukur dan yang tidak.
% Kasus Medikara: pemicu Medikara (insiden kebocoran data lewat vendor pada 2025, temuan audit
%   kontrol akses, antrean panjang, biaya integrasi tinggi) diturunkan ke strategi
%   Satu Pasien, Satu Rekam Medis.

\section{Model Operasi}

% TODO
% Cakupan: dua dimensi integrasi dan standardisasi proses menurut Ross, Weill, dan Robertson; empat
%   model operasi (diversifikasi, koordinasi, replikasi, unifikasi); konsekuensi tiap model pada
%   data dan aplikasi; core diagram sebagai gambaran satu halaman.
% Kasus Medikara: kondisi saat ini cenderung diversifikasi karena akuisisi; strategi baru menuntut
%   pergeseran ke koordinasi atau unifikasi.

\section{Fase Preliminary dan Architecture Vision}

% TODO
% Cakupan: kesiapan organisasi untuk kerja arsitektur; Request for Architecture Work; isi
%   Architecture Vision; business scenario singkat untuk menggali kebutuhan.

\section{Prinsip Arsitektur}

% TODO
% Cakupan: format prinsip TOGAF (nama, pernyataan, alasan, implikasi); ciri prinsip yang baik;
%   contoh prinsip yang buruk dan cara memperbaikinya; jumlah prinsip yang wajar.
% Kasus Medikara: prinsip satu sumber kebenaran untuk data pasien dan prinsip akses berdasarkan
%   kebutuhan.
% Pendalaman teknis: menurunkan prinsip menjadi aturan desain dan kriteria tinjauan rancangan yang
%   dapat diperiksa.

\section{Menulis Visi Arsitektur untuk Direksi}

% TODO
% Cakupan: struktur dokumen dua sampai tiga halaman; bahasa untuk pembaca non-teknis; nilai,
%   risiko, dan ruang lingkup yang eksplisit; cara meminta persetujuan.

\section{Aktivitas: Menyusun Visi Arsitektur}

% TODO
% Tujuan aktivitas: menghasilkan visi arsitektur yang menghubungkan strategi organisasi dengan
%   arah arsitektur.
% Peran dalam tim: analis bisnis memimpin telaah dokumen strategi dan wawancara; analis teknologi
%   menilai pemicu yang berasal dari sistem; seluruh tim menentukan model operasi dan prinsip;
%   koordinator merangkai dokumen visi arsitektur.
% Langkah: kumpulkan dokumen strategi; susun rantai motivasi dari pemicu sampai outcome terukur;
%   tentukan model operasi saat ini dan yang diinginkan; tulis 5 sampai 8 prinsip arsitektur;
%   rangkai visi arsitektur dua sampai tiga halaman.
% Luaran dan tonggak capstone: visi arsitektur dengan rantai motivasi, model operasi, dan prinsip
%   arsitektur, yang melengkapi C1.
% Berkas sources yang direncanakan: sources/chapter03/templat-rantai-motivasi.md,
%   sources/chapter03/lembar-penilaian-model-operasi.md,
%   sources/chapter03/templat-prinsip-arsitektur.md, sources/chapter03/templat-visi-arsitektur.md

\section{Ringkasan}

% TODO

\section{Latihan}

\begin{enumerate}
  \item Turunkan strategi ``Satu Pasien, Satu Rekam Medis'' menjadi satu \textit{goal}, dua \textit{outcome} terukur, dan satu \textit{requirement}.
  \item Jaringan 12 apotek menjalankan proses yang sama di semua cabang, tetapi pelanggan di tiap kota berbeda dan jarang berpindah cabang. Model operasi mana yang paling cocok, dan apa akibatnya bagi aplikasi dan data?
  \item Kritik prinsip ``Semua sistem harus memakai teknologi terbaru.'' Tulis ulang menjadi prinsip yang baik lengkap dengan alasan dan implikasi.
\end{enumerate}
